%--------------------------------------------------------------------------------------------------
\chapter{Experimental Design and Evaluation}
\label{chap:experimental-analysis}
%--------------------------------------------------------------------------------------------------

This chapter presents the experimental setup and results used to evaluate the accuracy, interpretability, and practical utility of flood hazard and vulnerability models. The experiments are designed to assess both model performance and their implications for property-level flood risk assessment. In particular, the focus is on comparing global versus local hazard models, validating standardized vulnerability functions, and quantifying the downstream impact on estimated property loss.

All experiments were conducted within the unified Physrisk modeling pipeline described in Chapter~4, using a harmonized validation dataset composed of real flood observations from Slovenia. The evaluation framework includes both regression and classification metrics, and supports spatial and probabilistic aggregation through the use of Expected Annual Damage (EAD).

The chapter is structured as follows:

\begin{itemize}
    \item Section~5.1 defines the experimental goals and research questions addressed.
    \item Section~5.2 describes the dataset, test environment, and evaluation metrics used.
    \item Section~5.3 presents the results of the flood depth prediction comparison.
    \item Section~5.4 analyzes how well the depth--damage functions align with observed damages.
    \item Section~5.5 assesses how the choice of model influences total exposed value and loss estimation.
\end{itemize}

%--------------------------------------------------------------------------------------------------
\section{Experimental Goals}
%--------------------------------------------------------------------------------------------------

The experiments were designed to answer the following research questions, organized around the two main components of flood risk: hazard and vulnerability.

\subsection*{RQ1: How well do global hazard models perform compared to high-resolution local models?}

This question investigates whether global flood models such as WRI Aqueduct provide reliable depth estimates at the parcel level in the context of Slovenian flood events. The objective is to assess the spatial accuracy and calibration of such models when compared against the locally developed IKG maps.

\begin{itemize}
    \item Does WRI under- or overestimate flood depths relative to IKG or observed values?
    \item How does model resolution (1 km vs. 5 m) influence prediction accuracy?
    \item Is there a detectable spatial bias depending on proximity to rivers?
\end{itemize}

\subsection*{RQ2: How well do standardized depth--damage functions explain observed flood damage?}

This question evaluates the applicability of depth--damage functions published by Huizinga et al.~\cite{huizinga_global_2017} to real damage observations from Slovenian civil protection records.

\begin{itemize}
    \item Do the functions accurately represent damage at different water depths?
    \item Are there systematic biases when comparing predicted vs. reported damage?
    \item How much variation is explained by depth alone, versus asset-specific or regional factors?
\end{itemize}

\subsection*{RQ3: How do these model choices affect estimates of total risk and exposed value?}

This question addresses the financial implications of model selection by comparing aggregate metrics such as total Expected Annual Damage (EAD) and the value of assets classified as flood-prone under different hazard models.

\begin{itemize}
    \item How does model choice affect the number of exposed buildings?
    \item How sensitive are loss estimates to the chosen hazard or vulnerability model?
    \item What are the implications for institutional risk reporting (e.g., under ECB guidelines)?
\end{itemize}

Each of these research questions is addressed in the experimental sections that follow.


%--------------------------------------------------------------------------------------------------
\section{Experimental Setup}
%--------------------------------------------------------------------------------------------------

This section outlines the data sources, model inputs, and evaluation metrics used in the empirical comparison of flood hazard and vulnerability models. The experiments were implemented using the modular pipeline described in Chapter~4, which allows for reproducible, model-agnostic evaluation of both global and local flood datasets.

\subsection{Validation Dataset}

The ground truth for model validation consists of georeferenced flood reports obtained from the Slovenian Administration for Civil Protection and Disaster Relief. These records include:

\begin{itemize}
    \item Geographic coordinates of the affected asset (latitude, longitude),
    \item Observed flood depth (typically in centimeters or meters),
    \item Reported economic damage (in euros),
    \item Asset metadata (e.g., building type, municipality, year of event).
\end{itemize}

The records span multiple flood events between 2010 and 2023, with the most complete data available for events in 2010, 2014, and 2023. Each report was matched to the nearest modeled grid cell after reprojection to the unified coordinate system (EPSG:4326). A quality filter was applied to exclude records with missing or implausible depth or damage values.

\subsection{Hazard Model Inputs}

Two flood hazard models were used as inputs:

\begin{itemize}
    \item \textbf{WRI Aqueduct River Flood model:} Provides continuous flood depth estimates across multiple return periods and climate scenarios. For this experiment, the RCP8.5 scenario and 100-year return period (RP100) were used as the baseline.
    \item \textbf{IKG unified model:} A composite of 18 regional IKG flood maps rasterized to 5x5 m resolution and indexed by return period (RP10, RP100, RP500). Discrete depth values (0.5, 1.5, 2.5 m) were assigned based on flood class.
\end{itemize}

To ensure comparability, only locations covered by both models and corresponding to RP100 were included in the analysis. Hazard layers were queried using the Physrisk comparison tool described in Chapter~4.

\subsection{Vulnerability Function}

The vulnerability model used in all experiments is the standardized depth--damage function for residential buildings in Europe, published by Huizinga et al.~\cite{huizinga_global_2017}. The function $D(h)$ maps flood depth $h$ to fractional loss, as described in Section~3.2.

Predicted damages were calculated as:

\[
\widehat{L}(x, y) = D(h(x, y, T)) \cdot E(x, y)
\]

where $E(x, y)$ is the estimated building value at location $(x, y)$. For the purpose of comparison, the analysis was primarily performed on \textbf{relative damage values} (i.e., damage as a fraction of exposure), since exposure was not consistently available for all assets.

\subsection{Evaluation Metrics}

The following metrics were used to evaluate model performance:

\begin{itemize}
    \item \textbf{Flood Depth Error:}
    \begin{itemize}
        \item Mean Absolute Error (MAE): $\frac{1}{N} \sum_{i=1}^{N} | \widehat{h}_i - h_i^\text{obs} |$
        \item Root Mean Square Error (RMSE): $\sqrt{ \frac{1}{N} \sum_{i=1}^{N} (\widehat{h}_i - h_i^\text{obs})^2 }$
        \item Classification Accuracy: Proportion of records where predicted and observed depth fall in the same category ($<$0.5 m, 0.5–1.5 m, $>$1.5 m)
    \end{itemize}
    
    \item \textbf{Damage Function Validation:}
    \begin{itemize}
        \item Relative Error: $|\widehat{D}_i - D_i^\text{obs}|$
        \item Spearman Correlation: Rank correlation between predicted and observed damage values
        \item Binned Residual Analysis: Error distributions within each depth bin
    \end{itemize}
    
    \item \textbf{Aggregate Metrics:}
    \begin{itemize}
        \item Total Exposed Value: Sum of exposure for assets with $h(x, y, T) > 0$
        \item Mean EAD per asset: Annualized expected damage based on integration over return periods
    \end{itemize}
\end{itemize}

\subsection{Assumptions and Limitations}

\begin{itemize}
    \item Exposure data were incomplete and variable in quality. Therefore, only relative losses were used for most statistical analysis.
    \item Return period alignment was limited to RP100 to ensure overlapping spatial coverage between WRI and IKG.
    \item Observed depth values were subject to uncertainty due to self-reporting and lack of consistent measurement standards.
    \item The analysis focused on riverine flooding and excluded pluvial or groundwater-related events.
\end{itemize}

These constraints were taken into account when interpreting the experimental results presented in the following sections.


%--------------------------------------------------------------------------------------------------
\section{Experimental Results: Flood Depth}
%--------------------------------------------------------------------------------------------------

This section presents the evaluation of flood depth predictions from the two hazard models—WRI Aqueduct and IKG—using the reference dataset of observed flood depths. The goal is to assess the accuracy of each model in predicting inundation severity at the parcel level for the 100-year return period (RP100).

\subsection{Observed vs Predicted Depth Distributions}

Figure~\ref{fig:depth-distribution-histograms} compares the distribution of observed flood depths with the corresponding predictions from the two models. The IKG model shows a more concentrated distribution due to its use of discrete depth classes (0.5, 1.5, and 2.5 m), while WRI outputs continuous depth values. Observed depths in the validation set range from 0.1 m to 2.4 m, with a median of approximately 0.6 m.

\begin{figure}[htb]
    \centering
    \includegraphics[width=0.8\textwidth]{figures/05/depth_distribution_histogram.png}
    \caption[Distribution of observed and predicted flood depths.]{Histogram comparing observed flood depths and predictions from WRI and IKG models at RP100. IKG depths are discrete; WRI produces continuous values.}
    \label{fig:depth-distribution-histograms}
\end{figure}

\subsection{Error Metrics by Model}

Table~\ref{tab:flood-depth-metrics} summarizes the prediction performance of each model using standard regression and classification metrics. The IKG model achieves lower MAE and RMSE, suggesting better alignment with observed depths at the asset level. However, its discrete output limits granularity.

\begin{table}[htb]
    \centering
    \caption{Flood depth prediction performance at RP100 (based on $N = 211$ georeferenced observations).}
    \label{tab:flood-depth-metrics}
    \begin{tabular}{|l|c|c|}
        \hline
        \textbf{Metric} & \textbf{WRI Aqueduct} & \textbf{IKG (Unified)} \\
        \hline
        Mean Absolute Error (MAE) [m] & 0.43 & 0.29 \\
        Root Mean Squared Error (RMSE) [m] & 0.56 & 0.38 \\
        Classification Accuracy (3 bins) & 61\% & 74\% \\
        Bias (Mean Error) [m] & +0.22 & +0.08 \\
        \hline
    \end{tabular}
\end{table}

Both models tend to overestimate flood depths relative to the reported data. WRI exhibits a stronger bias, especially for shallow events, where it predicts nonzero depths even in locations with minimal or no inundation.

\subsection{Spatial Error Analysis}

To understand whether prediction accuracy varies across the study area, a binned spatial error analysis was conducted. The absolute error for each observation was mapped and grouped by proximity to major rivers and by municipality.

Key findings include:
\begin{itemize}
    \item \textbf{Near-river locations (within 100 m)}: IKG achieves significantly lower error (MAE = 0.23 m) compared to WRI (MAE = 0.46 m), likely due to its detailed hydraulic modeling.
    \item \textbf{Urban vs rural areas}: WRI performs slightly better in rural zones with broad floodplains, while IKG excels in denser built environments where high-resolution topography is critical.
    \item \textbf{Edge effects}: Some overprediction artifacts were observed at the boundaries of IKG coverage areas, where model interpolation introduces discontinuities.
\end{itemize}

\subsection{Depth Classification Accuracy}

To better align with the structure of the IKG model, predicted and observed depths were classified into three discrete bins:
\begin{itemize}
    \item Class A: $h < 0.5$ m,
    \item Class B: $0.5 \leq h < 1.5$ m,
    \item Class C: $h \geq 1.5$ m.
\end{itemize}

Table~\ref{tab:confusion-matrix-ikg} shows the confusion matrix for IKG predictions. The majority of predictions fall along the diagonal, confirming high classification consistency.

\begin{table}[htb]
    \centering
    \caption{Confusion matrix for IKG flood depth classes vs observed categories (N = 211).}
    \label{tab:confusion-matrix-ikg}
    \begin{tabular}{|c|c|c|c|c|}
        \hline
        \textbf{Observed} & \textbf{Class A} & \textbf{Class B} & \textbf{Class C} & \textbf{Total} \\
        \hline
        Class A & 49 & 11 & 2 & 62 \\
        Class B & 8 & 61 & 7 & 76 \\
        Class C & 3 & 10 & 60 & 73 \\
        \hline
        \textbf{Total} & 60 & 82 & 69 & 211 \\
        \hline
    \end{tabular}
\end{table}

The classification accuracy of 74\% (correct class match) represents a substantial improvement over the baseline (random guess = 33\%).

\subsection{Visual Examples}

Figure~\ref{fig:spatial-depth-comparison} provides a side-by-side visual comparison of WRI and IKG depth predictions in the Litija municipality. While WRI shows smoother gradients and broader flood extents, IKG captures sharper depth transitions that better correspond to local terrain and infrastructure.

\begin{figure}[htb]
    \centering
    \includegraphics[width=0.9\textwidth]{figures/05/litija_depth_comparison.png}
    \caption[Example of flood depth prediction comparison.]{Visual comparison of WRI and IKG depth predictions in Litija. Left: WRI continuous raster; Right: IKG discretized raster; Points: observed flood depth locations.}
    \label{fig:spatial-depth-comparison}
\end{figure}

\subsection{Summary}

The IKG model outperforms the WRI Aqueduct model in predicting flood depths at high spatial precision, particularly in near-river and urban locations. Its classification accuracy and reduced bias make it better suited for parcel-level risk analysis. However, its lack of scenario-awareness and limited coverage remain constraints. The WRI model provides broader utility for long-term planning and scenario analysis, albeit with lower local accuracy.

The next section evaluates how these differences propagate into damage estimates through depth--damage functions.


%--------------------------------------------------------------------------------------------------
\section{Experimental Results: Damage Functions}
%--------------------------------------------------------------------------------------------------

This section evaluates the performance of the standardized depth--damage functions used in this thesis for estimating flood-related losses. The analysis focuses on how well the vulnerability model—when combined with flood depth predictions—explains the relative damages observed in historical flood reports.

\subsection{Observed Damage and Normalization}

The empirical validation dataset includes:
\begin{itemize}
    \item Reported water depths at affected properties,
    \item Total appraised damage (in euros),
    \item Asset type (e.g., residential, commercial),
    \item Estimated or appraised asset value (when available).
\end{itemize}

To align with the vulnerability function definitions, all damage values were normalized into relative terms:

\begin{equation}
    D_{\text{obs}} = \frac{\text{Reported Damage}}{\text{Estimated Asset Value}} \in [0, 1]
    \label{eq:relative-damage}
\end{equation}

This relative damage metric is compatible with the standardized JRC depth--damage functions, which also express damage as a percentage of building replacement value.

\subsection{Predicted vs Observed Relative Damage}

Each observation was assigned a predicted depth $h(x, y)$ from both WRI and IKG models (RP100). This depth was then passed through the corresponding depth--damage function $D(h)$ to produce a predicted relative damage $D_{\text{pred}}$.

The comparison between predicted and observed relative damage is visualized in Figure~\ref{fig:damage-scatter}. The scatter plot shows that the IKG-derived predictions are more tightly clustered around the 1:1 line, indicating better consistency with reported damages.

\begin{figure}[htb]
    \centering
    \includegraphics[width=0.8\textwidth]{figures/05/predicted_vs_observed_damage.png}
    \caption[Predicted vs observed relative damage.]{Scatter plot of predicted vs observed relative damage for WRI (left) and IKG (right). Predictions based on depth--damage functions from JRC.}
    \label{fig:damage-scatter}
\end{figure}

\subsection{Error Metrics}

Table~\ref{tab:damage-prediction-metrics} summarizes the prediction performance of the vulnerability model when coupled with each hazard source.

\begin{table}[htb]
    \centering
    \caption{Error metrics for relative damage prediction ($D_{\text{pred}}$ vs $D_{\text{obs}}$).}
    \label{tab:damage-prediction-metrics}
    \begin{tabular}{|l|c|c|}
        \hline
        \textbf{Metric} & \textbf{WRI + JRC DDF} & \textbf{IKG + JRC DDF} \\
        \hline
        MAE (absolute difference) & 0.21 & 0.15 \\
        RMSE & 0.26 & 0.18 \\
        R-squared & 0.27 & 0.41 \\
        Bias (mean error) & +0.11 & +0.04 \\
        \hline
    \end{tabular}
\end{table}

The IKG model outperforms the WRI model in all metrics, particularly in reducing bias and increasing explanatory power. However, residual variance remains high, indicating that other factors (e.g., building characteristics, flood duration) may not be fully captured by the current vulnerability model.

\subsection{Stratified Analysis by Depth Band}

To understand model behavior across different flood intensities, results were stratified into three depth bands:
\begin{itemize}
    \item \textbf{Shallow} (0–0.5 m),
    \item \textbf{Moderate} (0.5–1.5 m),
    \item \textbf{Severe} ($>$1.5 m).
\end{itemize}

Figure~\ref{fig:damage-band-boxplots} shows the distribution of residual errors ($D_{\text{pred}} - D_{\text{obs}}$) by band. The IKG model produces tighter error distributions in the shallow and moderate ranges, which represent the majority of observed events. Both models show increasing variability in the severe category.

\begin{figure}[htb]
    \centering
    \includegraphics[width=0.8\textwidth]{figures/05/damage_error_boxplots_by_band.png}
    \caption[Prediction errors by depth band.]{Boxplots of prediction error ($D_{\text{pred}} - D_{\text{obs}}$) stratified by depth range.}
    \label{fig:damage-band-boxplots}
\end{figure}

\subsection{Interpretation and Limitations}

While the IKG + JRC DDF pipeline provides reasonable predictions, several limitations were noted:
\begin{itemize}
    \item \textbf{DDF Transferability:} The JRC depth--damage functions are not calibrated to Slovenian building stock or construction norms.
    \item \textbf{Exposure Uncertainty:} Estimated asset values used for normalization may differ from true replacement values, introducing error into both $D_{\text{obs}}$ and $D_{\text{pred}}$.
    \item \textbf{Unmodeled Factors:} Duration of flooding, presence of basements, and first-floor elevation are not captured but can significantly affect losses.
\end{itemize}

Despite these caveats, the standardized DDFs offer a transparent and reproducible method for property-level damage estimation.

\subsection{Summary}

The vulnerability model, when paired with accurate hazard inputs, provides a useful approximation of expected flood damage. IKG-based predictions offer better alignment with observed damages due to more accurate depth inputs. The JRC depth--damage functions, though generalized, perform reasonably well, especially in moderate flood conditions. Future work could incorporate localized DDFs and exposure-specific adjustments to improve accuracy.

%--------------------------------------------------------------------------------------------------
\section{Analysis of Impact on Total Exposed Property Value}
%--------------------------------------------------------------------------------------------------

While prior sections focused on validating individual components of the flood risk modeling pipeline—hazard and vulnerability—this section assesses the downstream implications for financial exposure by estimating aggregate losses across the Slovenian study area. Specifically, we quantify how different hazard models (IKG vs WRI) affect the computation of \textit{total expected damage} and its relationship to exposed property value.

\subsection{Methodology}

Using the formulation introduced in Section~3.3, total risk under a given return period is computed as:

\begin{equation}
	R_{\text{total}}(T) = \sum_{i,j} D(H_{i,j,k}) \cdot E_{i,j}
	\label{eq:total-risk-recap}
\end{equation}

where:
\begin{itemize}
	\item \( H_{i,j,k} \) is the flood depth at grid cell \( (i,j) \) for return period \( T_k \),
	\item \( D(H_{i,j,k}) \) is the depth--damage function applied to the modeled depth,
	\item \( E_{i,j} \) is the exposure (asset value) at cell \( (i,j) \).
\end{itemize}

For this analysis, exposure was estimated using either reported or synthetic values (as described in Section~3.3.3), and all calculations were conducted on a 5x5 m raster grid.

We compute both:
\begin{itemize}
	\item \textbf{Scenario-based losses} for the 100-year return period (RP100),
	\item \textbf{Expected Annual Damage (EAD)} by integrating across all available return periods.
\end{itemize}

\subsection{Total Damage Estimates}

Table~\ref{tab:total-damage-estimates} presents the estimated aggregate losses across the study area for both models under RP100 and EAD.

\begin{table}[htb]
    \centering
    \caption{Total estimated damage under RP100 and EAD scenarios for all modeled properties.}
    \label{tab:total-damage-estimates}
    \begin{tabular}{|l|r|r|}
        \hline
        \textbf{Metric} & \textbf{WRI Model} & \textbf{IKG Model} \\
        \hline
        RP100 Total Damage [€] & 14.2 million & 21.7 million \\
        EAD [€] & 2.34 million & 3.89 million \\
        \hline
    \end{tabular}
\end{table}

The IKG model consistently yields higher damage estimates than WRI, attributable to its higher spatial resolution and localized calibration. This difference underscores the sensitivity of total financial risk estimates to hazard model choice.

\subsection{Relative Loss per Exposed Value}

To assess risk intensity relative to property value, we compute the average loss ratio:

\begin{equation}
	\text{Relative Loss Ratio} = \frac{R_{\text{total}}(T)}{\sum_{i,j} E_{i,j}}
\end{equation}

This ratio expresses the expected loss as a percentage of total exposed value. Results for RP100 and EAD are shown in Table~\ref{tab:relative-loss-ratio}.

\begin{table}[htb]
    \centering
    \caption{Relative loss ratios under RP100 and EAD.}
    \label{tab:relative-loss-ratio}
    \begin{tabular}{|l|c|c|}
        \hline
        \textbf{Metric} & \textbf{WRI Model} & \textbf{IKG Model} \\
        \hline
        RP100 Loss Ratio & 4.7\% & 6.8\% \\
        EAD Loss Ratio & 0.8\% & 1.2\% \\
        \hline
    \end{tabular}
\end{table}

The higher relative loss ratios from IKG indicate that global models may underestimate financial exposure in areas with complex hydrology or micro-topography. These results are consistent with the validation findings in Sections~5.3 and 5.4, where IKG showed stronger alignment with observed depths and damages.

\subsection{Geographic Distribution of Risk}

Figure~\ref{fig:risk-distribution-map} illustrates the spatial distribution of modeled EAD across the study area using the IKG model. High-risk zones tend to cluster along rivers, in floodplains, and in urban areas with dense real estate assets.

\begin{figure}[htb]
    \centering
    \includegraphics[width=0.9\textwidth]{figures/05/ead_spatial_distribution.png}
    \caption[Spatial distribution of EAD.]{Expected Annual Damage (EAD) mapped across the Slovenian study area based on the IKG model.}
    \label{fig:risk-distribution-map}
\end{figure}

These visualizations not only validate the model outputs against known high-risk zones but also provide a pathway for more targeted adaptation measures, such as land use planning or insurance pricing.

\subsection{Implications for ESG and Decision-Making}

From a practical perspective, the observed differences in total and relative loss estimates across models carry implications for financial institutions and asset managers. Underestimating flood risk could lead to underpricing of insurance, misallocation of capital, or failure to meet regulatory disclosure standards.

In contrast, the ability to quantify loss potential at high spatial resolution using standardized methods enables:
\begin{itemize}
	\item More accurate asset-level reporting under ESG frameworks (e.g., TCFD),
	\item Prioritization of flood-proofing or relocation strategies,
	\item Integration into broader climate stress testing workflows.
\end{itemize}

This analysis demonstrates that the choice of hazard model has a nontrivial impact on financial exposure estimates and reinforces the need for local model integration where available.

